\begin{abstract}
Short-term solar photovoltaic (PV) prediction is difficult when plant output is shaped by both deterministic solar geometry and fast-changing cloud conditions. This paper evaluates a physics-guided two-stage XGBoost workflow for 5-min-grid AC power prediction at a \PlantCapacity{} solar PV plant located at \StudySite. The reported experiment is an offline target-weather-input evaluation: target-time meteorological and cloud columns are taken from the held-out dataset, while deployment requires those fields to be supplied before issue time by a forecast service or local weather pipeline. The first model learns a clearness-index-like transmissivity signal from meteorological, cloud-layer, calendar, and solar-geometry features. A second XGBoost model then combines this learned atmospheric signal with clear-sky and operational features to estimate plant power, followed by explicit clipping to the plant and theoretical-power envelope. The evaluation uses a chronological train-validation-test protocol over a \DataSpan{} operational plant dataset containing 525,888 records after machine-fault, hardware-error, shutdown, and invalid-sensor periods were omitted during data cleaning. On the held-out daylight test period, the proposed two-stage path obtains an RMSE of \RMSEproposed, an nRMSE of \NRMSEproposed, an MAE of \MAEproposed, and an $R^2$ of \RtwoProposed. These values measure performance under the fixed target-weather-input protocol, not under archived issue-time weather forecasts at specific lead times. The model improves over previous-day same-time persistence but does not beat a live 5-min last-value nowcast baseline; reporting both baselines separates forecast-input evaluation from sensor-nowcast evaluation. Ablation results regenerated with the active pipeline show a measurable gain over direct single-stage XGBoost and a smaller gain over a total-cloud-only cascade. The diagnostics examine weather-regime errors, month-hour behavior, calibration, ramp-rate sensitivity, feature importance, baseline skill, and daily energy tracking. The contribution is a traceable plant-level prediction workflow and deployment template rather than a new machine-learning algorithm.
\end{abstract}

\begin{keywords}
Clearness index, cloud cover, photovoltaic power forecasting, physics-guided machine learning, solar energy, XGBoost.
\end{keywords}
